% Transfert direct d'électrons : lame de zinc dans une solution de sulfate de cuivre(II) (L'essentiel, chapitre 8)
\begin{tikzpicture}
\pic at (0,0) {becher={largeur=2.2,hauteur=2.2,niveau=1.6,liquide=pcBleu}};
\filldraw[fill=pcZinc,draw=otInk!80,line width=0.5pt] (-0.15,0.32) rectangle (0.15,2.85);
\pic at (6,0) {becher={largeur=2.2,hauteur=2.2,niveau=1.6,liquide=pcBleu!40}};
\filldraw[fill=pcZinc,draw=otInk!80,line width=0.5pt] (5.85,0.32) rectangle (6.15,2.85);
\filldraw[fill=pcCuivre,draw=pcCuivre!70!black,line width=0.3pt] (5.8,0.32) rectangle (6.2,1.6);
\foreach \yy in {0.45,0.78,1.11,1.44}{%
  \fill[pcCuivre!75!black] (5.78,\yy) circle (0.05) (6.22,{\yy+0.14}) circle (0.05);}
\draw[pcFleche] (1.6,1.25) -- (4.4,1.25);
\node[pcLegende,anchor=south] at (3,1.35) {quelques minutes};
\pcNote[west]{(0.15,2.45)}{(1.35,2.8)}{lame de zinc}
\pcNote[west]{(6.22,1.0)}{(7.45,1.3)}{dépôt de cuivre}
\node[pcLegende,anchor=north,align=center] at (0,-0.15) {solution de sulfate de cuivre(II)};
\node[pcLegende,anchor=north,align=center] at (6,-0.15) {la solution pâlit :\\les ions \ce{Cu^2+} disparaissent};
\end{tikzpicture}
\par\smallskip\legende{Les ions cuivre(II) captent directement les électrons cédés par le zinc : du cuivre se dépose et des ions zinc(II) passent en solution.}
